\section{When the contrapositive reveals a usable representation}
Now reverse the statement: if $n^2$ is even, must $n$ be even? The equation $n^2=2k$ is less convenient because it does not directly express $n$ as twice an integer. A different approach is to suppose $n$ is not even.

Every integer is either even or odd, but not both. This follows from division by two with remainder $0$ or $1$, a basic integer-arithmetic fact. Thus the contrapositive says: if $n$ is odd, then $n^2$ is odd. Write $n=2k+1$ and calculate
\begin{align*}
n^2&=(2k+1)(2k+1)\\
&=4k^2+4k+1\\
&=2(2k^2+2k)+1.
\end{align*}
The final expression is odd. Therefore $n^2$ cannot be even when $n$ is odd. Combining this with the first direction proves
\[
n\text{ is even}\quad\Longleftrightarrow\quad n^2\text{ is even}.
\]

The strategic lesson is specific: the negation of the desired conclusion supplied a more useful representation than the original hypothesis did. Contraposition is not a ritual to apply whenever a proof feels difficult.

\subsection*{What a contradiction permits, and what it does not}
To prove a claim by contradiction, temporarily suppose the claim is false. Retain the original hypotheses. If their consequences force a statement and its denial, the temporary supposition cannot hold. We then conclude the claim is true. This is ordinary classical reasoning, the setting of this book.

For example, suppose an integer were both even and odd. Then $2k=2\ell+1$ for integers $k,\ell$, and subtraction would give $2(k-\ell)=1$. Division gives $k-\ell=1/2$, impossible because the left side is an integer. The contradiction is that the same number would be an integer and equal to a noninteger. Saying merely ``that looks unlikely'' would not finish the proof.

In the next argument the contradiction is not a false numerical equality. It is a conflict with a deliberately chosen \emph{smallest} denominator. We will explain why that choice exists before using it. A contradiction can violate a selection property just as decisively as it can violate arithmetic.
